\documentclass[11pt]{article}
\usepackage[margin=1in]{geometry}
\usepackage[T1]{fontenc}
\usepackage{lmodern,microtype,hyperref,enumitem,amsmath,amssymb,xcolor,booktabs}
\usepackage[numbers,sort&compress]{natbib}
\usepackage{url}
\hypersetup{colorlinks=true,urlcolor=blue,citecolor=blue,linkcolor=blue}

\title{Q13 --- Comprehension check}
\author{DLS 2026 --- AI Internals Module}
\date{}

\begin{document}
\maketitle

\section*{Knowledge representation, storage, and retrieval in transformers}

A central puzzle in the study of large language models is \emph{where} factual
knowledge physically lives inside the network, and \emph{how} it is surfaced at
inference time. The dominant empirical account points to the feed-forward
(MLP) sublayers as key--value memory stores: the first weight matrix acts as a
bank of pattern-detecting ``keys'' that fire on specific input contexts, and
the second matrix contributes corresponding ``value'' vectors that are written
into the residual stream \citep{geva2021kvmemories}. Causal interventions
reinforce this picture: \citet{meng2022rome} show that a small set of
mid-layer MLP modules are decisive for recalling a given subject--relation fact,
and that editing those weights (ROME) can surgically alter the stored
association. Complementary work attributes facts to neurons that, when ablated
or amplified, change the predicted object \citep{dai2022knowledgeneurons}.
Attention heads then play a retrieval role, moving information between token
positions---induction heads, for instance, implement a copy-and-complete
mechanism that underlies in-context recall \citep{olsson2022induction}. The
interpretability toolkit that connects these components treats the residual
stream as a shared communication channel written and read by each layer
\citep{elhage2021mathematical}, with intermediate ``logit lens'' readouts
revealing how a prediction is progressively refined
\citep{geva2022promoting}.

\paragraph{Question.}
A model reliably completes the prompt ``\emph{The Eiffel Tower is located in
the city of}'' with ``\emph{Paris}.'' (a) Using the key--value memory and
causal-tracing perspectives, describe where this fact is most likely stored
and by which mechanism it is retrieved and written to the output. (b) An
engineer applies ROME to a mid-layer MLP so that the model now answers
``\emph{Rome}.'' Explain concretely what ROME changes and why editing that
specific component suffices, and name one limitation such weight edits exhibit.

\paragraph{Answer.}
(a) The fact is stored predominantly in the second (output) weight matrix of
one or more \emph{mid-layer} MLP sublayers, operating as key--value memory
\citep{geva2021kvmemories}. As the subject token ``Eiffel Tower'' is processed,
attention heads first aggregate the relevant subject information at the last
token position---a routing step consistent with the induction/copy mechanisms
and residual-stream communication framework
\citep{olsson2022induction,elhage2021mathematical}. At a mid-layer MLP, a
subset of keys in the first matrix detects this subject-plus-relation pattern
(``located-in-city'' context), and the paired value vectors in the second
matrix are added into the residual stream, incrementing the component of the
representation aligned with the ``Paris'' output direction. Causal tracing
localizes this effect: corrupting the subject tokens and then restoring
individual hidden states shows that restoring the mid-layer MLP outputs at the
last subject token recovers the correct prediction, identifying those modules
as the decisive site of storage \citep{meng2022rome}. The ``logit lens'' view
shows the ``Paris'' prediction becoming readable in the residual stream after
these layers write to it \citep{geva2022promoting}, and the final unembedding
projects this to a high logit for ``Paris.''

(b) ROME treats the second MLP matrix $W$ as a linear associative memory and
performs a rank-one update $W' = W + \Lambda (C^{-1} k_*)^{\top}$ that inserts a
new key--value association: it computes the key vector $k_*$ that the subject
``Eiffel Tower'' induces at the targeted layer, and chooses $\Lambda$ so the
retrieved value $W' k_*$ now points toward the ``Rome'' output direction while
minimally disturbing other stored associations (hence the $C^{-1}$
covariance-whitening term) \citep{meng2022rome}. Editing this one mid-layer MLP
suffices precisely because causal tracing identified it as the pivotal
read/write site for subject--relation facts; changing where the value is
written at that bottleneck propagates to the final logits. A key limitation is
that such weight edits generalize poorly and can damage unrelated knowledge:
they may fail to update logically entailed or paraphrased queries, can be
undone or confused by related edits, and risk collateral drift in nearby
associations---so the edit is local and brittle rather than a robust belief
update \citep{meng2022rome,dai2022knowledgeneurons}.

\bibliographystyle{plainnat}
\begin{thebibliography}{9}

\bibitem[Dai et al.(2022)]{dai2022knowledgeneurons}
Damai Dai, Li Dong, Yaru Hao, Zhifang Sui, Baobao Chang, and Furu Wei.
\newblock Knowledge neurons in pretrained transformers.
\newblock In \emph{Proceedings of the 60th Annual Meeting of the Association
for Computational Linguistics (ACL)}, 2022.

\bibitem[Elhage et al.(2021)]{elhage2021mathematical}
Nelson Elhage, Neel Nanda, Catherine Olsson, et al.
\newblock A mathematical framework for transformer circuits.
\newblock \emph{Transformer Circuits Thread}, 2021.
\newblock \url{https://transformer-circuits.pub/2021/framework/index.html}.

\bibitem[Geva et al.(2021)]{geva2021kvmemories}
Mor Geva, Roei Schuster, Jonathan Berant, and Omer Levy.
\newblock Transformer feed-forward layers are key-value memories.
\newblock In \emph{Proceedings of the 2021 Conference on Empirical Methods in
Natural Language Processing (EMNLP)}, 2021.

\bibitem[Geva et al.(2022)]{geva2022promoting}
Mor Geva, Avi Caciularu, Kevin Ro Wang, and Yoav Goldberg.
\newblock Transformer feed-forward layers build predictions by promoting
concepts in the vocabulary space.
\newblock In \emph{Proceedings of the 2022 Conference on Empirical Methods in
Natural Language Processing (EMNLP)}, 2022.

\bibitem[Meng et al.(2022)]{meng2022rome}
Kevin Meng, David Bau, Alex Andonian, and Yonatan Belinkov.
\newblock Locating and editing factual associations in GPT.
\newblock In \emph{Advances in Neural Information Processing Systems (NeurIPS)},
2022.

\bibitem[Olsson et al.(2022)]{olsson2022induction}
Catherine Olsson, Nelson Elhage, Neel Nanda, et al.
\newblock In-context learning and induction heads.
\newblock \emph{Transformer Circuits Thread}, 2022.
\newblock \url{https://transformer-circuits.pub/2022/in-context-learning-and-induction-heads/index.html}.

\end{thebibliography}

\end{document}